% !TeX root = zk-frame-gauge.tex
\documentclass[11pt]{article}
\usepackage[T1]{fontenc}
\usepackage{lmodern,microtype}
\usepackage[margin=1in]{geometry}
\usepackage{amsmath,amssymb,amsthm}
\usepackage[colorlinks=true,allcolors=blue]{hyperref}
% =====================================================================
%  Notation.  One definition per notion, for the whole document.
%
%  A symbol means the same thing in the main matter and in both annexes.
%  Where the source notes were merged, the annexes were adapted to the
%  main matter, and where a letter was already taken the annex object was
%  given a free symbol through the semantic macros below.  The map back to
%  the notation of the cited papers is in Annex \ref{annex:pcs}, "Symbols".
%
%  Multilinear extensions are written \widetilde{\cdot} throughout
%  (\mle below).  The hat is reserved for the normalized subspace
%  polynomials of the additive NTT (\Wn).
% =====================================================================

% ---------- fields ----------
\newcommand{\F}{\mathbb{F}}
\newcommand{\Ftwo}{\mathbb{F}_2}
\newcommand{\K}{K}            % F_{2^64}:  addresses, counters, committed lanes
\newcommand{\E}{E}            % F_{2^192} = K[y]/(y^3+y+1): words and challenges
\newcommand{\gen}{\mathsf{g}} % generator of K^*, order 2^64 - 1

% ---------- checked relations ----------
\newcommand{\qeq}{\stackrel{?}{=}}  % an equality the verifier checks, not one that holds by construction

% ---------- multilinears ----------
\newcommand{\mle}[1]{\widetilde{#1}}
\newcommand{\eq}{\operatorname{eq}}
\newcommand{\cube}[1]{\{0,1\}^{#1}}
\newcommand{\inner}[2]{\langle #1,\, #2 \rangle}
\newcommand{\fc}{\chi}      % a sumcheck round challenge, everywhere in the document
\newcommand{\flock}{\mathrm{flock}}

% ---------- Flock protocol (Annex C) ----------
\newcommand{\qflock}{q_{\flock}}             % packed Boolean Flock witness
\newcommand{\qext}[1]{#1^{\sharp}}           % 64-point univariate / multilinear extension
\newcommand{\kblock}{\kappa_{\mathrm{block}}} % variables in one BLAKE2s witness block
\newcommand{\kbatch}{\kappa_{\mathrm{batch}}} % log number of BLAKE2s instances
\newcommand{\kbool}{\kappa_{\mathrm{bool}}}   % variables of the unpacked Boolean batch witness
\newcommand{\kskip}{\kappa_{\mathrm{skip}}}   % variables replaced by the univariate skip
\newcommand{\skipdom}{\mathcal{H}}            % 64 skip-domain nodes
\newcommand{\skipcoset}{\mathcal{H}'}         % 64 transmitted round-one nodes
\newcommand{\fembed}{\varphi_{8}}             % embedding of F_{2^8} into E
\newcommand{\zskip}{\upsilon_{\mathrm{zc}}}   % zerocheck univariate challenge
\newcommand{\lcalpha}{\alpha_{\mathrm{lc}}}      % lincheck batching challenge (A, B, C, the pin)
\newcommand{\ksk}{k_{\mathrm{sk}}}              % row index, the skipped half
\newcommand{\kin}{k_{\mathrm{in}}}              % row index, the within-block half
\newcommand{\jsk}{j_{\mathrm{sk}}}              % column index, the skipped half
\newcommand{\jin}{j_{\mathrm{in}}}              % column index, the within-block half
\newcommand{\jconst}{j_{\mathbf{1}}}            % position of the circuit's constant-one wire
\newcommand{\posof}[1]{k(#1)}                 % position, equivalently R1CS row, of a committed wire
\newcommand{\lform}[1]{\mathcal{F}_{#1}}        % a wire's expansion, over the block's positions
\newcommand{\erow}{e_{\mathrm{row}}}                % lincheck's row weight, the quirky eq table
\newcommand{\wcol}{w_{\mathrm{col}}}                % lincheck's column weight
\newcommand{\chg}[1]{\Gamma_{#1}}                 % a wire's coefficient in the backward walk
\newcommand{\colmar}{M_{\mathrm{col}}}           % column marginal of one matrix, A_0 or B_0

% ---------- VM ----------
\newcommand{\loc}[1]{\mem[\fp\cdot #1]}
\newcommand{\dmode}{\mathsf{mode}}   % a DEREF store mode
\newcommand{\pc}{\mathbf{pc}}
\newcommand{\fp}{\mathbf{fp}}
\newcommand{\nxt}[1]{\mathrm{next}(#1)}   % the successor of a register: next(pc), next(fp)
\newcommand{\mem}{\mathbf{m}}
\newcommand{\addr}{\mathit{addr}}
\newcommand{\val}{\mathit{val}}
\newcommand{\cnt}{\mathit{count}}
\newcommand{\md}{\mathrm{md}}       % BLAKE2s metadata: byte counter, final and last-node flags
\newcommand{\cntfin}{\mathit{count}^{\mathrm{fin}}}  % committed finalize-count column
\newcommand{\rcnts}{\mathcal{N}}          % multiset of read counts at one (address, value)
\newcommand{\mset}[1]{\{\!\{#1\}\!\}}     % a multiset (plain braces stay sets)
\newcommand{\lut}{\mathbf{L}}             % a lookup array (memory, or the program)
\newcommand{\sep}{\dsep{sep}}             % the domain separator of a generic lookup array
\newcommand{\ncomp}{N_{\mathrm{comp}}}    % components of a lookup entry (its width)
\newcommand{\push}{\textsc{push}}
\newcommand{\pull}{\textsc{pull}}
\newcommand{\tup}[1]{( #1 )}
\newcommand{\rdf}[1]{\dsep{read}\,\tup{#1}}   % a pull/push pair differing only in the count
\newcommand{\dsep}[1]{\mathsf{#1}}
\newcommand{\opc}[1]{\mathsf{op}_{\mathsf{#1}}}
\newcommand{\srcsel}[1]{\mathsf{src}(#1)}
\newcommand{\nprog}{N_{\mathrm{prog}}}   % program length
\newcommand{\nmem}{N_{\mathrm{mem}}}    % memory length
\newcommand{\kbc}{\kappa_{\mathrm{bc}}}   % log program length: nprog = 2^kbc 
\newcommand{\kmem}{\kappa_{\mathrm{mem}}}  % log memory length: nmem = 2^kmem
\newcommand{\nread}{N_{\mathrm{read}}}   % total read flushes in a proof
\newcommand{\ntab}{N_{\mathrm{tab}}}     % number of tables (fixed by the arithmetization)
\newcommand{\nside}{N_{\mathrm{side}}}   % grand-product sides: push, pull, count
\newcommand{\taumax}{\tau_{\max}}      % rounds of the table sumcheck: the tallest table's log height
\newcommand{\ncol}[1]{N_{\mathrm{col},#1}}  % columns of table #1
\newcommand{\ncons}[1]{N_{\mathrm{cons},#1}}  % constraints of table #1
\newcommand{\nconstot}{N_{\mathrm{cons}}}    % constraints of all tables, the batching offset
\newcommand{\nfl}[1]{N_{\mathrm{flushes},#1}}    % bus flushes per row of table #1
\newcommand{\btup}{\mathbf{f}}          % a bus tuple
\newcommand{\bpull}{\mathsf{pull}}        % the state a row pulls from the bus
\newcommand{\bpush}{\mathsf{push}}        % the state a row pushes onto the bus

% ---------- codes and the PCS (Annex B) ----------
% Letters that the main matter already spends on something else are routed
% through these, so the annex reads with one meaning per symbol.
\newcommand{\Enc}{\mathrm{Enc}}
\newcommand{\kdim}{k}                 % message dimension, the convention
\newcommand{\RS}{\mathrm{RS}}
\newcommand{\nv}{M}                      % variable count of a level   (was m_i)
\newcommand{\nlev}{\Lambda}              % number of levels            (was r)
\newcommand{\rate}{\rho}                 % code rate, the convention
\newcommand{\rexp}{\nu}                  % rate exponent, rate = 2^-nu  (was R)
\newcommand{\rad}{\gamma}                % proximity radius, the convention
\newcommand{\slack}{\eta}                % Johnson slack, the convention
\newcommand{\mcapar}{\theta}            % MCA parameter, the convention
\newcommand{\dmin}{\delta_{\min}}        % relative minimum distance   (was delta)
\newcommand{\rc}[1]{c^{(#1)}}            % round-#1 running claim      (was sigma_j)
\newcommand{\nq}{\omega}                 % query count                 (was t_i)
\newcommand{\qset}{\mathcal{Q}}          % sampled column indices      (was J_i)
\newcommand{\blen}{n}                    % code block length, the convention
\newcommand{\dom}{\mathcal{D}}           % evaluation domain           (was S)
\newcommand{\orc}{Z}                     % oracle / interleaved word   (was C)
\newcommand{\nrows}{N_{\mathrm{rows}}}    % committed rows of the level-0 oracle (R is the bus root)
\newcommand{\nstack}{N_{\mathrm{stack}}}  % field elements placed in the stack, before the zero padding
\newcommand{\lst}{\mathcal{L}}           % list of nearby codewords    (was Lambda)
\newcommand{\ffinal}{f^{\ast}}           % final plaintext table       (was p)
\newcommand{\cood}{c_{\mathrm{ood}}}     % out-of-domain answer        (was y)
\newcommand{\mcanum}{A_{\mathrm{mca}}}   % MCA error numerator         (was a)
\newcommand{\mcanumi}[1]{A_{\mathrm{mca},#1}}  % ... at level #1
\newcommand{\polyset}{\mathcal{G}}       % polynomials bound by a commitment
\newcommand{\relopen}{\mathsf{R}_{\mathrm{open}}}
\newcommand{\subsp}{\mathcal{U}}         % F_2-subspace                (was U_i)
\newcommand{\vansub}{\mathcal{V}}        % subspace vanishing poly     (was W_i)
\newcommand{\Wn}[1]{\widehat{\vansub}_{#1}} % its normalization (hat = normalized)
\newcommand{\novel}{\mathcal{X}}         % novel-basis polynomial      (was X_j)
\newcommand{\ntmap}{\psi}                % linearized NTT map          (was q_d)
\newcommand{\bfarr}{\mathsf{B}}          % NTT butterfly array         (was B)
\newcommand{\mpoly}{\mathcal{P}}         % message polynomial          (was P_u)
\newcommand{\aset}{\mathcal{A}}          % Johnson agreement set       (was A_j)

% ---------- ring switching (Annex A) ----------
\newcommand{\pair}{\Omega}               % error pairing values        (was V_k)
\newcommand{\supp}{\mathcal{S}}          % support of \Phimap          (was S)
\newcommand{\nflock}{\kappa_{\flock}}     % variables of the packed flock witness (was m, L)
\newcommand{\lag}{\varpi}                 % Lagrange weights of flock's skip (was p_i)
\newcommand{\fcoord}{\mathsf{a}}          % an F_2 coordinate function  (was A_w)

% ---------- statistical privacy research ----------
\newcommand{\zkSec}{\kappa_{\mathrm{zk}}}
\newcommand{\zkStmt}{\mathsf{x}}
\newcommand{\zkWit}{\mathsf{w}}
\newcommand{\zkAux}{\mathsf{z}}
\newcommand{\zkRel}{\mathsf{R}_{\mathrm{VM}}}
\newcommand{\zkProver}{\mathsf{P}}
\newcommand{\zkVerifier}{\mathsf{V}}
\newcommand{\zkSim}{\mathsf{Sim}}
\newcommand{\zkView}{\operatorname{View}}
\newcommand{\zkSD}{\operatorname{SD}}
\newcommand{\zkErr}{\varepsilon_{\mathrm{zk}}}
\newcommand{\zkBits}{b_{\mathrm{zk}}}
\newcommand{\zkHybrids}{m_{\mathrm{zk}}}
\newcommand{\zkLaneLen}{H_{\mathrm{lane}}}
\newcommand{\zkLaneLog}{\ell_{\mathrm{lane}}}
\newcommand{\zkPadBlock}{P_{\mathrm{pad}}}
\newcommand{\zkFrameSize}{s_{\mathrm{frame}}}
\newcommand{\zkFreeLanes}{J_{\mathrm{free}}}
\newcommand{\zkPinSpace}{W_{\mathrm{pin}}}
\newcommand{\zkPrefixSpace}{T_{\mathrm{pin}}}
\newcommand{\zkMemoryMasks}{D_{\mathrm{mem}}}
\newcommand{\zkProbePrefix}{L_{\mathrm{probe}}}
\newcommand{\zkProbeLog}{\ell_{\mathrm{probe}}}
\newcommand{\zkProbeCountSize}{L_{\mathrm{count}}}
\newcommand{\zkProbeCountQuota}{T_{\mathrm{probe}}}
\newcommand{\zkProbeAccess}{A_{\mathrm{probe}}}
\newcommand{\zkProbeCompleted}{C_{\mathrm{probe}}}
\newcommand{\zkProbeReadBudget}{B_{\mathrm{probe}}}
\newcommand{\zkProbeRepeat}{R_{\mathrm{probe}}}
\newcommand{\zkLeafPublicProbes}{B_{\mathrm{leaf}}}
\newcommand{\zkLeafDigitProbes}{D_{\mathrm{leaf}}}
\newcommand{\zkLeafDigitCap}{t_{\mathrm{digit}}}
\newcommand{\zkLeafXmss}{N_{\mathrm{leaf,X}}}
\newcommand{\zkLeafSphincs}{N_{\mathrm{leaf,S}}}
\newcommand{\zkLeafWords}{N_{\mathrm{word}}}
\newcommand{\zkDigitPolynomial}{P_{\mathrm{digit}}}
\newcommand{\zkPatchBcCenter}[1]{A^{\mathrm{bc}}_{\mathrm{patch},#1}}
\newcommand{\zkPatchMemCenter}{A^{\mathrm{mem}}_{\mathrm{patch}}}
\newcommand{\zkPatchRouter}[1]{R_{\mathrm{patch},#1}}
\newcommand{\zkCompactSupport}{S_{2,\mathrm{c}}}
\newcommand{\zkCompactError}{\delta_{2,\mathrm{c}}}
\newcommand{\zkPlacedSize}{N_{\mathrm{placed}}}
\newcommand{\zkMemoryAux}{\mathcal A_{\mathrm{mem}}}
\newcommand{\zkInitialImage}{\mathcal I_0}
\newcommand{\zkLeafError}{\varepsilon_{\mathrm{leaf}}}
\newcommand{\zkFlockOrder}{\pi_{\mathrm F}}
\newcommand{\zkProfilePoly}[1]{d_{#1}}
\newcommand{\zkCosetMinor}{\mathcal D_{\mathrm{cos}}}
\newcommand{\zkCosetOffsets}{S_{\mathrm{cos}}}
\newcommand{\zkPatternCount}{G_{\mathrm{pat}}}
\newcommand{\zkRealFactor}{T_{\mathrm{real}}}
\newcommand{\zkQueryMasks}{M_{\mathrm{query}}}
\newcommand{\zkQueryShift}{d_{\mathrm{query}}}
\newcommand{\zkPinnedCV}{H_{\mathrm{pin}}}
\newcommand{\zkMetadataTest}{\ell_{\mathrm{meta}}}
\newcommand{\zkMetadataBank}{B_{\mathrm{meta}}}
\newcommand{\zkRealRows}{I_{\mathrm{real}}}
\newcommand{\zkTripleSupport}{S_{3,\mathrm{c}}}
\newcommand{\zkTripleError}{\delta_{3,\mathrm{c}}}
\newcommand{\zkBlockCollapse}[1]{\mathcal C_{#1}}
\newcommand{\zkMatchingError}{\delta_{\mathrm{match}}}
\newcommand{\zkFixedError}[1]{\delta_{\mathrm{fixed}}(#1)}
\newcommand{\zkFixedMap}{\mathcal L_{\mathrm{fixed}}}
\newcommand{\zkSwapCoins}{\mathsf{bits}}
\newcommand{\zkDenseSupport}{S_{\mathrm{dense}}}
\newcommand{\zkDenseError}[1]{\delta_{\mathrm{dense}}(#1)}
\newcommand{\zkLeafCounts}{\mathbf c_{\mathrm{leaf}}}
\newcommand{\zkSharedPcSupport}{S_{\mathrm{pc}}}
\newcommand{\zkLeafResidual}{R_{\mathrm{leaf}}}
\newcommand{\zkShortSupport}{S_{11}}
\newcommand{\zkShortError}{\delta_{11}}
\newcommand{\zkChildPcSupport}{S_{\mathrm{pc},4}}
\newcommand{\zkGkrChildren}{\mathbf h_{\mathrm{GKR}}}
\newcommand{\zkChildPoint}{\zeta^{\circ}}
\newcommand{\zkCosetLow}[1]{H_{#1}^{\mathrm{cos}}}
\newcommand{\zkWideCosetLow}[1]{H_{#1}^{(16)}}
\newcommand{\zkWideCosetInvariant}[1]{I_{#1}^{(16)}}
\newcommand{\zkWideCosetNoise}{\mathcal M_{16}}
\newcommand{\zkCosetDual}[1]{D_{#1}^{(16)}}
\newcommand{\zkWideCosetError}{\varepsilon_{\mathrm{cos},12}}
\newcommand{\zkJointQueryMap}{A_{\mathrm{jq}}}
\newcommand{\zkPublicQueryMap}{A_{\mathrm{pub}}}
\newcommand{\zkPublicQueries}{J_{\mathrm{pub}}}
\newcommand{\zkJointQueryDefect}{\delta_{\mathrm{jq}}}
\newcommand{\zkJointCoins}{\theta_{\mathrm{jq}}}
\newcommand{\zkPublicQueryShift}{b_{\mathrm{pub}}}
\newcommand{\zkSixQuerySet}{J_6}
\newcommand{\zkSixQueryTest}{T_6}
\newcommand{\zkLowPairs}{\mathcal B_{\mathrm{low}}}
\newcommand{\zkLowWeights}{a^{\mathrm{low}}}
\newcommand{\zkPreQueryCoins}{\theta_{\mathrm{pre}}}
\newcommand{\zkJointNoiseCode}{\mathcal M_{\mathrm{jq}}}
\newcommand{\zkJointShiftSpace}{\mathcal T_{\mathrm{jq}}}
\newcommand{\zkDualSupports}{\mathcal S_{\mathrm{jq}}}
\newcommand{\zkDualSupportCount}[1]{N_{#1}^{\mathrm{jq}}}
\newcommand{\zkThirtyTwoLow}[1]{H_{#1}^{(32)}}
\newcommand{\zkThirtyTwoLower}[2]{L_{#1,#2}^{(32)}}
\newcommand{\zkThirtyTwoUpper}[2]{U_{#1,#2}^{(32)}}
\newcommand{\zkBalancedLowWeight}[1]{a_{#1}^{(32)}}
\newcommand{\zkBalancedHighWeight}[1]{b_{#1}^{(32)}}
\newcommand{\zkLowQueryBand}{\mathcal A_{13}}
\newcommand{\zkLowBandNoiseCode}{\mathcal M_{\mathrm{low},13}}
\newcommand{\zkScatteredQueries}{J_{42}}
\newcommand{\zkCountFrontier}[1]{F^{\mathrm{cnt}}_{14,#1}}
\newcommand{\zkCountHeadChild}[1]{H^{\mathrm{cnt}}_{14,#1}}
\newcommand{\zkCountCompletionFactor}{A^{\mathrm{cnt}}_{14,0}}
\newcommand{\zkFineCountFrontier}[1]{F^{\mathrm{cnt}}_{4,#1}}
\newcommand{\zkFineCountChild}{H^{\mathrm{cnt}}_{4,0}}
\newcommand{\zkAdapterError}{\delta_{\mathrm{adapter}}}
\newcommand{\zkCalibrationCap}{C_{\mathrm{cal}}}
\newcommand{\zkCalibrationLarge}{A_{\mathrm{cal}}}
\newcommand{\zkCalibrationSmall}{B_{\mathrm{cal}}}
\newcommand{\zkCalibrationTarget}{S_{\mathrm{cal}}}
\newcommand{\zkBcRouterHalf}{h_{\mathrm{bc}}}
\newcommand{\zkBcCoarseTarget}[1]{T^{\mathrm{bc}}_{#1}}
\newcommand{\zkBlakeLoad}{d_{\mathrm B}}
\newcommand{\zkMemoryCoarseCap}{C_{\mathrm{read}}}
\newcommand{\zkMemoryCoarseTarget}{T_{\mathrm{read}}}
\newcommand{\zkCopyRepeats}{R_{\mathrm{copy}}}
\newcommand{\zkRouterWidth}{k_{\mathrm{route}}}
\newcommand{\zkRouterRadius}{A_{\mathrm{route}}}
\newcommand{\zkRouterLength}{N_{\mathrm{route}}}
\newcommand{\zkAllocationCap}{s_{\mathrm{alloc}}}

\newtheorem{lemma}{Lemma}
\title{Address-preserving randomness for control products}
\author{}
\date{}
\begin{document}
\maketitle
\vspace{-2em}

\paragraph{Result and scope.} Public random frame scalings and private cycle choices can jointly hide six JUMP products while retaining eight correlated fields. This supplies control randomness absent from the previous private payload masks. The construction uses the optional honest public setup and the documented field-valued ISA, but requires raw field-valued bytecode and padding-row support beyond the current native prover. It is a restricted statistical decoupling theorem under honest independent verifier coins, not a full GKR simulator or an implemented ZK mode.

\section{Where payload randomness has no effect}

In the depth-$22$ layout of \path{zk-gkr-second-wire.tex}, the second push packet's child $1$ is
\[
 \eq(x,0)P_1+\eq(x,1)P_5+\eq(x,2)P_9+\eq(x,3)P_{13}.
\]
Here $P_1$ is memory quarter $1$'s seed product, $P_5$ is public bytecode quarter $1$, $P_9$ is MUL bytecode, and $P_{13}$ multiplies the JUMP state and bytecode blocks. Fix the coins, public SET library and private occupancy choices. The previous MUL-by-one and XOR-by-zero payloads occupy memory quarter $0$, while unused private words occupy quarter $2$; none changes these four products. Increasing those payload banks alone therefore cannot randomize this scalar. The audit replays both layers through the actual verifier using products of valid local cycles, with all omitted factors set to one. It is not a complete VM trace or an attack on the remaining occupancy randomness.

\paragraph{Frame scaling.} At SET address $p$, choose a physical frame $h\in\K^\times$ and a public scale $a\in\K^\times$. Use
\[
 \fp=ha,\qquad o_{\rm SET}=a^{-1},\qquad
 (o_c,o_d,o_f)=(\gen,\gen^2,\gen^3)/a.
\]
The closing JUMP at $\gen p$ reads condition $1$, destination $p$ and frame value $ha$. Every accessed address remains one of $h,h\gen,h\gen^2,h\gen^3$. Install both alternatives' bytecode and memory values before choosing either privately, then execute the chosen SET/JUMP cycle twice. Images are fixed across private choices, all count leaves are unchanged, and the state multiset is balanced. The scales are independent uniform public values; private fair occupancy bits supply the hiding entropy.

\paragraph{Representation boundary.} The ISA uses $\K$-valued frames and operands and bounds accessed addresses. Its equations permit this scaling. In contrast, \path{crates/lean_vm/src/cpu/isa.rs} and \path{trace.rs} represent offsets and frame indices as \texttt{u32}. For example $a=\gen^{2^{40}}$ gives frame exponent $2^{18}+2^{40}$ and SET offset exponent $2^{64}-1-2^{40}$, although their product accesses $\gen^{2^{18}}$. Both exponents exceed the native types. No verifier equation or weight is added here, but native trace construction is not yet supported, and no implementation cost claim follows.

\section{Jointly treat the control fingerprints}

Write $Q=|\K|$, $q=|\E|=Q^3$, $e_i=\eq(\alpha,i)$ and $t_i=\alpha_i/(1+\alpha_i)$. Exclude $\alpha_i\in\{0,1\}$ and put $z=\beta/e_0$. Conditional on $\alpha$, $z$ is uniform in $\E$. Let $\sigma_{\rm s},\sigma_{\rm c},\sigma_{\rm m}$ be the distinct state, bytecode and memory tags, and let $\omega=\gen^{\texttt{OP\_JUMP}}$. Define
\[
\begin{aligned}
 s_0(a)&=e_0(z+\sigma_{\rm s}+t_0p+t_1ha),&
 s_1(a)&=s_0(a)+e_0t_0(\gen+1)p,\\
 B&=z+\sigma_{\rm c}+t_0\gen p+t_0t_1\omega,&
 D&=\gen+t_0\gen^2+t_1\gen^3,\\
 X_\ell(a)&=e_0(B+t_1\gen^\ell+t_2D/a),&
 L_\ell(a)&=e_0(z+\sigma_{\rm m}+t_0h\gen^3+t_1\gen^\ell+t_0t_1ha).
\end{aligned}
\]
The two executions contribute, in state, bytecode and frame-memory push/pull order,
\[
 U(a)=\bigl(s_0^2,s_1^2,X_1X_2,X_0X_1,L_1L_2,L_0L_1\bigr).
\]
The normalized roots in $a$ comprise two state roots, three reciprocal bytecode roots and three frame roots:
\[
 \xi_0=\frac{z+\sigma_{\rm s}+t_0p}{t_1h},\quad
 \xi_1=\frac{z+\sigma_{\rm s}+t_0\gen p}{t_1h},\quad
 \zeta_\ell=\frac{t_2D}{B+t_1\gen^\ell},\quad
 \nu_\ell=\frac{z+\sigma_{\rm m}+t_0h\gen^3+t_1\gen^\ell}{t_0t_1h}.
\]

\begin{lemma}[Regular control factors]
Suppose $D$ and all three bytecode denominators are nonzero, all eight roots have degree three over $\K$, and their Frobenius orbits are pairwise disjoint. For uniform $a\in\K^\times$, every nontrivial character of $U(a)\in(\E^\times)^6$ has mean of magnitude at most $\kappa=24\sqrt Q/(Q-1)$.
\end{lemma}
\begin{proof}
Adjoin the root $0$ to handle the reciprocal factors. The eight roots and zero generate the finite \'etale $\K$-algebra $\E^8\times\K$, of dimension $25$. A character with exponents $(c_0,\ldots,c_5)$ induces root exponents
\[
 (2c_0,2c_1,c_3,c_2+c_3,c_2,c_5,c_4+c_5,c_4,-2c_2-2c_3).
\]
This character remains nontrivial: multiplication by $2$ is invertible modulo the odd order $q-1$, and the other four exponents appear directly. Leading coefficients only contribute phases. Katz's finite \'etale character-sum theorem gives $24\sqrt Q$ for the sum over $a\in\K$, extending by zero at $a=0$~\cite{katz}. Divide by $Q-1$. All other factors are nonzero on $\K$ by regularity.
\end{proof}

\section{Avoid a base-field-sized exceptional probability}

Consider $N$ private choice bits, with distinct SET and closing-JUMP addresses across their $2N$ alternatives, and distinct physical frames. Outside a coin event of probability at most
\[
 \delta=\frac{8+6N}{q}+\frac{1+158N}{q-2}
              +\frac{7Q^2}{q(q-2)},
\]
discarding at most seven bits makes every remaining alternative regular. Discarding here means conditioning on those bits in the analysis, not deleting trace rows. The discarded indices depend on the coins and fixed addresses, not on the public scales.

\paragraph{Small-field roots.} A state root belongs to $\K$ precisely when $u+rp'\in\K$, where $r=t_0/t_1$, $u=(z+\sigma_{\rm s})/t_1$, and $p'$ runs over distinct instruction addresses. Unless both $r,u\in\K$, at most one address qualifies: two would force $r\in\K$. The joint exception costs at most $Q^2/[q(q-2)]$, because $r$ has point masses at most $1/(q-2)$ and $u$ is independently uniform. For each $\ell$, bytecode roots satisfy the same test with slope $t_0/(t_2D)$ and variable address $\gen p$; frame roots use slope $1/t_1$ and address $h\gen^3$, after multiplying the root by $h\in\K^\times$. Each family has the same joint exception bound. There are $1+3+3$ families, hence at most seven bad addresses, and therefore seven bad bits. Excluding $D=0$ costs $1/(q-2)$; the $6N$ zero bytecode denominators cost $6N/q$.

\paragraph{Within-family conjugacy.} The two state roots differ by $t_0(\gen+1)p/(t_1h)$; a frame-root pair differs by $(\gen^\ell+\gen^m)/(t_0h)$. In each case one root is uniform independently of the ratio in that difference, whose point masses are at most $1/(q-2)$. Each nontrivial Frobenius equality therefore costs at most $1/(q-2)$. This gives two comparisons for states and six for frame roots.

For a bytecode pair, put $b_\ell=B+t_1\gen^\ell$. The equation $(t_2D/b_\ell)^{Q^j}=t_2D/b_m$, $j=1,2$, has at most $Q-1$ solutions in $t_2$. It is soluble only if $\operatorname{Norm}(b_\ell)=\operatorname{Norm}(b_m)$. As $z$ varies uniformly, $b_\ell/b_m$ ranges bijectively over $\E\setminus\{1\}$, apart from the pole. Exactly $Q^2+Q$ of these ratios have norm one. Thus the probability is at most $(Q^2+Q)(Q-1)/[q(q-2)]<1/(q-2)$. There are six comparisons; unconjugated roots within each three-root family are distinct.

\paragraph{Cross-family conjugacy.} State and frame roots are independent of $t_2$, while each bytecode root is a nonzero multiple of $t_2$. Each equality with any of three Frobenius conjugates fixes at most one $t_2$, giving $18+27$ comparisons per alternative. For a state root $x$ at address $p'\in\{p,\gen p\}$, set
\[
 C=\sigma_{\rm m}+\sigma_{\rm s}+t_0(h\gen^3+p').
 \qquad
 \nu_\ell=\frac{x}{t_0}+\frac{C}{t_0t_1h}+\frac{\gen^\ell}{t_0h}.
\]
The distinct tags ensure that $C=0$ excludes at most one $t_0$ per state address. After excluding it, each equation $x^{Q^j}=\nu_\ell$ fixes at most one $t_1$ conditional on $(x,t_0)$; $x$ is uniform independently of these challenges. This gives $18$ further comparisons and two degeneracy exclusions per alternative. Altogether the conjugacy and degeneracy cost is $2N(2+6+6+18+27+18+2)/(q-2)=158N/(q-2)$. These union bounds are taken before conditioning on the other good events.

\section{A six-product statistical reduction}

Use the four banks with $n=3840$ bits each from \path{zk-set-payload-products.tex}, so $N=4n$. Fix their SET immediate words before sampling independent public scales. Retain eight actual fields: the JUMP condition-memory and destination-memory push/pull products, the SET payload-memory push/pull products, and the memory and bytecode final-counter evaluations. For fixed coins and private bits, none depends on the scales. Their joint law is retained, not simulated.

Two independent alternatives give a ratio character bound $b=\kappa^2<2^{-54}$. Apply the nonlinear-side-information lemma of \path{zk-public-seed-leakage.tex} with $G=(\E^\times)^6$, $m=q^8$, $c=(1+b)/2$, and conservatively $n-7$ usable bits in every bank. In this intermediate hybrid replace zero factors in discarded bits by one; the final zero bound pays for that change. Conditioning on discarded bits and their scales then only changes nonzero fixed cofactors and the retained map; averaging them back preserves the bound. The mean conditional distance satisfies
\[
 \eta^2\le\frac{q^8}{4}
 \left[\bigl(1+((q-1)^6-1)c^{3833}\bigr)^4-1\right]
 <2^{-1142}.
\]
Multiplying each of the six coordinates across banks preserves its independent uniform nonzero target. Averaging over good honest coins and applying Markov, the distance contribution is below $2^{-160}$ outside a fraction below $2^{-411}$ of honestly sampled scale libraries. Public SET words are fixed independently of those scales; choosing them after seeing the scales would invalidate the retained-map hypothesis.

For each fixed library, zeros among the alternatives' $16N$ state, bytecode and frame-memory factors, and at most $M\le2^{24}$ complementary factors, cost at most $(M+16N)/q$ over uniform $\beta$. This is charged separately, not used to condition setup on a witness. Consequently the six full JUMP products decouple from the actual retained eight-field law with error
\[
 2^{-160}+\delta+\frac{M+16N}{q}<2^{-159}.
\]
Nonzero multiplication by fixed disjoint complementary products and addition of fixed evaluation contributions preserve the bound. Those contributions may still depend on the witness. The setup exception concerns the fixed padding family and retained map, not a witness-dependent selection of good libraries.

\paragraph{Composition boundary.} The six products determine the banks' contribution to both node-$13$ products, and the SET state products are the reversed JUMP state products. SET bytecode products, memory finalization and later messages depend jointly on scales and private choices; they cannot be retained by merely counting their output bits. The continuation \path{zk-control-composition.tex} supplies an ordered public-word/scale hybrid, reducing the enlarged residual to seven fields and handling the second child-$1$ pair after removing the two state products. Other second-layer messages remain outside that theorem; a complete first/second-layer simulator is still required.

\paragraph{Evidence and size.} \path{zk_frame_gauge_audit.py} checks exact small-field root and conjugacy counts, actual-field formulas for all six products, valid cycles and unchanged count leaves, retained and excluded fields, the second-layer payload blind spot, and rational bounds. The $30720$ SET and closing-JUMP rows fit the previous heights $15$ and $17$, with memory and bytecode logs $20$ and stack log $24$. Scales add no rows or physical addresses to that reservation; they do change the public bytecode and padding memory. No full trace or production proof is generated.

\begin{thebibliography}{1}
\bibitem{katz} N.~M.~Katz. \emph{An Estimate for Character Sums}, Theorem~2. Journal of the American Mathematical Society 2(2), 1989, pp.~197 to 200. \href{https://web.math.princeton.edu/~nmk/old/estcharsums.pdf}{Author's copy}.
\end{thebibliography}
\end{document}
